% ==============================================================================
% CHƯƠNG 18: DỰ ÁN THỰC CHIẾN FRONTEND TỪ A-Z (FULL FRONTEND PROJECT MASTERCLASS)
% ==============================================================================
\chapter{Dự Án Thực Chiến Frontend Từ A-Z}
\label{chap:full_frontend_project}

Chương này hướng dẫn xây dựng trọn vẹn một ứng dụng Web thực tế: \term{Dashboard Quản Lý Khóa Học Trực Tuyến (EduWeb Admin Dashboard)} từ khâu phân tích bố cục, dựng HTML5 ngữ nghĩa, thiết kế CSS3 nâng cao đến xử lý logic JavaScript tương tác và Dark Mode.

% ==============================================================================
% 1. PHÂN TÍCH YÊU CẦU & THIẾT KẾ BỐ CỤC DỰ ÁN
% ==============================================================================
\section{Phân Tích Yêu Cầu \& Thiết Kế Bố Cục Dự Án}

\subsection{1. Cấu Trúc Thư Mục Dự Án Chuẩn Thực Tế}

\begin{codeWindow}[Cấu Trúc Thư Mục Dự Án Frontend Chuyên Nghiệp]
\begin{lstlisting}[language=HTTP]
eduweb-dashboard/
├── index.html              # Trang chu giao dien Dashboard
├── assets/
│   ├── css/
│   │   ├── main.css        # CSS Design System & Variables
│   │   └── dashboard.css   # Dynamic UI styles
│   ├── js/
│   │   ├── app.js          # Main Application Entry Point
│   │   └── courses.js      # Data Manager & Event Handlers
│   └── images/
│       └── logo.svg
└── README.md
\end{lstlisting}
\end{codeWindow}

\subsection{2. Bố Cục Giao Diện Kiến Trúc Grid \& Flexbox}

\begin{conceptBox}[Kiến Trúc Giao Diện 2 Cột (Sidebar + Main Area)]
\begin{itemize}[leftmargin=1.5em, itemsep=3pt]
    \item \textbf{Thanh Bên (Sidebar <aside>)}: Chứa Logo thương hiệu, Menu điều hướng chính và Nút đổi chế độ Sáng/Tối (\term{Dark Mode Toggle}).
    \item \textbf{Khu Vực Chính (<main>)}: Chứa Header tìm kiếm, Bộ lọc khóa học theo danh mục và Lưới sản phẩm (\term{Course Card Grid}).
\end{itemize}
\end{conceptBox}

% ==============================================================================
% 2. TRIỂN KHAI MÃ NGUỒN HTML, CSS & JAVASCRIPT
% ==============================================================================
\section{Triển Khai Mã Nguồn HTML, CSS \& JavaScript}

\subsection{1. Mã Nguồn Cấu Trúc HTML5 Masterclass}

\begin{codeWindow}[Mã Nguồn HTML5: Khung Xương Giao Diện Dashboard]
\begin{lstlisting}[language=HTML5]
<div class="dashboard-wrapper">
  (*@<!-- Thanh (*@bên@*) Sidebar -->@*)
  <aside class="sidebar">
    <div class="logo">EduWeb Admin</div>
    <nav class="sidebar-nav">
      <a href="#" class="active">(*@Khóa *@Học@*)</a>
      <a href="#">(*@Học *@Viên@*)</a>
      <a href="#">(*@Báo *@Cáo@*)</a>
    </nav>
    <button id="dark-mode-btn">(*@Đổi@*) Giao (*@Diện@*)</button>
  </aside>

  (*@<!-- Khu (*@vực nội@*) dung (*@chính@*) -->@*)
  <main class="main-content">
    <header class="top-bar">
      <search>
        <input type="search" id="search-input" placeholder="(*@Tìm kiếm khóa học...@*)">
      </search>
    </header>

    <section class="course-grid" id="course-list">
      <!-- Dynamic render by JS -->
    </section>
  </main>
</div>
\end{lstlisting}
\end{codeWindow}

\begin{browserWindow}[Kết Quả Hiển Thị Trình Duyệt: Mô Phỏng Giao Diện Dashboard Tổng Thể]
\sffamily\small
\begin{tcolorbox}[colback=gray!4, colframe=primaryBlue!30, arc=6pt, top=4pt, bottom=4pt, left=6pt, right=6pt]
\begin{tabularx}{\linewidth}{K{3.8cm} Z}
\begin{tcolorbox}[colback=primaryBlue, colframe=primaryBlue, arc=4pt, left=6pt, right=6pt, top=6pt, bottom=6pt]
\color{white}
\textbf{EduWeb Admin}\\[8pt]
\tcbox[on line, arc=3pt, colback=white!20, colframe=white!20, left=4pt, right=4pt, top=2pt, bottom=2pt]{\color{white}\bfseries $\bullet$ Khóa Học}\\[4pt]
$\circ$ Học Viên\\[3pt]
$\circ$ Báo Cáo\\[12pt]
\tcbox[on line, arc=3pt, colback=accentOrange, colframe=accentOrange, left=6pt, right=6pt, top=2pt, bottom=2pt]{\color{white}\tiny Dark Mode: Off}
\end{tcolorbox}
&
\begin{tcolorbox}[colback=white, colframe=gray!20, arc=4pt, left=8pt, right=8pt, top=6pt, bottom=6pt]
\textbf{Tìm kiếm:} \tcbox[on line, arc=3pt, colback=gray!8, colframe=gray!30, left=6pt, right=40pt, top=2pt, bottom=2pt]{\color{gray!70}Nhập tên khóa học...}\\[8pt]
\textbf{\color{primaryBlue} Danh Sách Khóa Học Nổi Bật:}\\[6pt]
\begin{tabular}{@{}l l@{}}
\begin{tcolorbox}[colback=white, colframe=primaryBlue!30, width=5.2cm, arc=4pt]
\textbf{\color{primaryBlue} HTML5 Masterclass 2026}\\[2pt]
\tiny Tác giả: Nguyễn Trọng Tấn\\[2pt]
\tcbox[on line, arc=2pt, colback=green!10, colframe=green!50, boxsep=1pt]{\color{green!60!black}\tiny Đang Hoạt Động}
\end{tcolorbox}
&
\begin{tcolorbox}[colback=white, colframe=primaryBlue!30, width=5.2cm, arc=4pt]
\textbf{\color{primaryBlue} Modern CSS Grid \& Flexbox}\\[2pt]
\tiny Tác giả: Đội Ngũ EduWeb\\[2pt]
\tcbox[on line, arc=2pt, colback=green!10, colframe=green!50, boxsep=1pt]{\color{green!60!black}\tiny Đang Hoạt Động}
\end{tcolorbox}
\end{tabular}
\end{tcolorbox}
\end{tabularx}
\end{tcolorbox}

\vspace{4pt}
\textbf{$\rightarrow$ Giải thích kết quả}: Bố cục 2 cột hoàn chỉnh với Sidebar bên trái chứa nút đổi giao diện và Main Area bên phải chứa thanh tìm kiếm cùng các Card khóa học render động.
\end{browserWindow}

\subsection{2. Triển Khai Logic JavaScript Tương Tác (Search \& Dark Mode)}

\begin{codeWindow}[Mã Nguồn JavaScript: Tìm Kiếm Thời Gian Thực \& Chế Độ Dark Mode]
\begin{lstlisting}[language=JavaScript]
// 1. Tinh nang tim kiem khoa hoc nhanh (Real-time Filter)
const searchInput = document.getElementById('search-input');
const courseCards = document.querySelectorAll('.course-card');

searchInput.addEventListener('input', function(e) {
  const searchTerm = e.target.value.toLowerCase().trim();

  courseCards.forEach(card => {
    const titleText = card.querySelector('.course-title').textContent.toLowerCase();
    if (titleText.includes(searchTerm {
      card.style.display = 'block'; (*@// Hien thi card trung (*@khớp@*)
    } else {
      card.style.display = 'none';  # An card khong (*@khớp@*)
    }
  });
});

// 2. Tinh nang chuyen doi Dark Mode luu vao localStorage
const darkModeBtn = document.getElementById('dark-mode-btn');
darkModeBtn.addEventListener('click', () => {
  document.body.classList.toggle('dark-theme');
  const isDark = document.body.classList.contains('dark-theme');
  localStorage.setItem('theme', isDark ? 'dark' : 'light');
});
\end{lstlisting}
\end{codeWindow}

\begin{browserWindow}[Kết Quả Hiển Thị Trình Duyệt: Mô Phỏng Khi Tìm Kiếm Từ Khóa "HTML5"]
\sffamily\small
\begin{tcolorbox}[colback=gray!90!black, colframe=gray!70!black, arc=6pt, top=4pt, bottom=4pt, left=6pt, right=6pt]
\color{white}
\textbf{Trạng Thái Dark Mode: ON} \hfill \textbf{Ô tìm kiếm:} \tcbox[on line, arc=3pt, colback=gray!80, colframe=gray!60, left=6pt, right=30pt, top=2pt, bottom=2pt]{\color{yellow}\bfseries HTML5}\\[6pt]
\begin{tcolorbox}[colback=gray!80!black, colframe=yellow!70!orange, width=0.85\linewidth, arc=4pt]
\color{white}
\textbf{\color{yellow} HTML5 Masterclass 2026} \hfill \tcbox[on line, arc=2pt, colback=yellow!20, colframe=yellow!80, boxsep=1pt]{\color{yellow!90!black}\tiny Match: 1 Kết Quả}\\[2pt]
{\small Tác giả: Nguyễn Trọng Tấn --- Khóa học lập trình HTML5 ngữ nghĩa từ căn bản đến nâng cao.}
\end{tcolorbox}
\end{tcolorbox}

\vspace{4pt}
\textbf{$\rightarrow$ Giải thích kết quả}: Khi gõ từ khóa "HTML5", JavaScript ngay lập tức lọc các card trùng khớp và giao diện tự động bật tone màu tối mượt mà nhờ Dark Mode!
\end{browserWindow}

% ==============================================================================
% 3. TÓM TẮT DỰ ÁN
% ==============================================================================
\section{Tóm Tắt \& Checklist Đóng Gói Đưa Lên Production}

\begin{itemize}[leftmargin=1.5em, itemsep=4pt]
    \item Đảm bảo mã HTML5 passing kiểm thử chuẩn W3C Validator không có lỗi đóng thẻ.
    \item Tối ưu dung lượng hình ảnh WebP/SVG và bật lazy loading cho các ảnh thumbnail.
    \item Kiểm thử giao diện Responsive trên màn hình di động 320px đến màn hình 4K.
    \item Kiểm tra điểm số Lighthouse đạt trên 90 điểm cho 4 tiêu chí: Performance, Accessibility, Best Practices và SEO.
\end{itemize}

Dự án thực chiến này đã tổng kết thành công toàn bộ hành trình học tập từ Nền tảng HTML5, CSS3 đến JavaScript Modern.
